\scanpage{200}
\mspage{74}
d'un objet comme $\underline{\lim}$ filtrante grande devant $\unsure{\pi|R}$ d'objets de $\mathrm{Filt}^\pi$). Je ne vais pas refaire le travail, fait \deleted{dans une \unsure{toute} \unsure{bonn}} \inserted{dans un cas} beaucoup plus général et plus délicat dans SGA 4 I §9, savoir dans le cor. 9.24 du th. \unsure{9.23|9.27|9.22}, (p. 177). La seule restriction, c'est qu'on \inserted{y} suppose les objets de $M$ \deleted{$\unsure{\pi|R}$-}accessibles, et désignant par $\pi_1$ \deleted{un} le plus petit cardinal tel que $\pi_1\geq\pi_0$ et que les objets de $\mathrm{Filt}^{\pi_0}$ soient $\pi_1$-accessibles, on se borne \inserted{\unsure{pour} la filtration cardinale de $\underline{\mathrm{Hom}}\,(S,M)$} aux cardinaux $\boxed{\pi\geq c=2^{c_0},}$ où
\[
 c_0=\mathrm{Sup}\,(\pi_1,\operatorname{Card}\mathrm{Fl}\, S).
\]
C'est \unsure{une} restriction \unsure{raisonnable} et \unsure{modérée}.

Pour le \unsure{'instant|l'instant}, je suis intéressé surtout au cas où $S=\Delta^1$, donc :
\[
 \underline{\mathrm{Hom}}\,(\Delta^1,M)=\underline{\mathrm{Fl}}\,(M),
\]
auquel cas on a $c_0=\pi_1$, donc il suffit de se borner aux cardinaux $\pi\geq 2^{\pi_1}$ pour avoir une bonne filtration cardinale sur $\underline{\mathrm{Fl}}\,(M)$.

C'est donc \unsure{l'ens|l'une} des deux conditions b), c) sur $M$ (p 44) qui est stable par passage aux $\underline{\mathrm{Hom}}\,(S,M)=M(S)$. Il en est de $=$ de la condition a) (stabilité pour petites limites ind \deleted{\unsure{filtrantes}}), et \inserted{aussi} de la condi%
% archive sheet number: 199
